\section*{Projektüberblick}
\addcontentsline{toc}{section}{Projektüberblick}

MIRA ist ein Softwareprototyp, der Glas, Metall, Papier, Kunststoff und Restmüll in Kamerabildern erkennen und lokalisieren soll. Untersucht wird, welcher messbare Zielkonflikt zwischen Erkennungsqualität, Modellgröße und CPU-Latenz entsteht, wenn ein kompaktes Detektionsmodell trainiert und auf INT8 quantisiert wird. Ein erster, abgeschlossener Machbarkeitsnachweis (Stage~A) verglich Bildklassifikatoren für vier Wertstoffklassen. Historische Detektionsläufe EXP-013 bis EXP-017 verwendeten anschließend YOLO11n und unterschiedliche Mischungen öffentlicher Datensätze; ihre Ergebnisse dienen als Ausgangspunkt, stammen aber nicht aus dem neu aufgebauten Datensatz. Dieser aktuelle, annähernd nach Objektannotationen ausgeglichene Fünfklassen-Datensatz umfasst getrennte Trainings-, Validierungs- und Testsätze. Der zugehörige Trainings- und Quantisierungsvergleich steht noch aus. Ein Raspberry Pi, eine Kamera und ein Sortierarm sind als spätere Integrationsstufe geplant und wurden in der vorliegenden Softwareuntersuchung noch nicht vermessen.
